\subsection{Phân tích các bên liên quan}

Việc phát triển và triển khai hệ thống mang lại giá trị gia tăng rõ rệt cho các bên liên quan thông qua việc giải quyết trực tiếp các nhu cầu thực tế trong vận hành và kinh doanh. Những nhu cầu này chính là cơ sở để xác lập danh mục các yêu cầu chức năng và phi chức năng chi tiết ở các mục tiếp theo.

\subsubsection{Khách hàng}

Đối với người dùng cuối, hệ thống cung cấp một nền tảng mua sắm hiện đại, tập trung vào tính tiện dụng và cá nhân hóa. Nhu cầu cốt lõi bắt đầu từ việc thiết lập tài khoản an toàn \textit{[FR1, FR2, FR3]} và quản lý thông tin cá nhân thuận tiện \textit{[FR4]}. Trong quá trình mua sắm, khách hàng cần các công cụ tìm kiếm, lọc sản phẩm minh bạch \textit{[FR5, FR8]} và quy trình giỏ hàng, thanh toán linh hoạt \textit{[FR6, FR7]}.

Điểm nhấn quan trọng là sự hỗ trợ thông minh từ trợ lý AI từ khâu tìm kiếm, gợi ý sản phẩm đến tư vấn các chính sách dịch vụ \textit{[FR10, FR11]}. Để thỏa mãn nhóm người dùng này, hệ thống phải đáp ứng các tiêu chuẩn về tính ổn định của giao diện \textit{[NFR1]}, tính dễ học trong thao tác \textit{[NFR3]} cùng khả năng phản hồi tự nhiên, đúng ngữ cảnh của AI Agent \textit{[NFR2]}.

\subsubsection{Đội ngũ vận hành và quản trị viên}

Đối với đội ngũ nội bộ, hệ thống đóng vai trò là công cụ tối ưu hóa hiệu suất làm việc hàng ngày. Nhu cầu chính là khả năng kiểm soát quyền hạn truy cập \textit{[FR14]} và quản lý danh mục sản phẩm một cách tập trung \textit{[FR15]}. Nhân viên vận hành cần được hỗ trợ cập nhật trạng thái đơn hàng chính xác \textit{[FR17]} và theo dõi các đánh giá từ khách hàng sau khi hoàn tất mua sắm \textit{[FR13]}.

Về phía kỹ thuật, nhóm này yêu cầu hệ thống phải duy trì tính toàn vẹn dữ liệu cao giữa các dịch vụ \textit{[NFR6]} và có khả năng bảo trì, nâng cấp linh hoạt nhờ vào kiến trúc phân tán \textit{[NFR5]}. Việc theo dõi và phát hiện các dịch vụ đang hoạt động giúp đội ngũ vận hành duy trì độ tin cậy của toàn bộ hệ thống \textit{[NFR7]}.

\subsubsection{Chủ sở hữu doanh nghiệp}

Về mặt chiến lược kinh doanh, hệ thống được kỳ vọng sẽ nâng cao hiệu quả tiếp cận thị trường và tối ưu hóa lợi nhuận thông qua các báo cáo thống kê trực quan về doanh thu và sản phẩm \textit{[FR16]}. Việc cho phép khách hàng chủ động theo dõi lịch sử đơn hàng \textit{[FR12]} và thực hiện quyền hủy đơn theo ràng buộc \textit{[FR9]} giúp xây dựng niềm tin thương hiệu và sự chuyên nghiệp cho doanh nghiệp.

Ngoài ra, hệ thống phải đáp ứng đòi hỏi khắt khe về an toàn tài sản dữ liệu thông qua các cơ chế bảo mật nghiêm ngặt \textit{[NFR4]}. Kiến trúc hệ thống cần đảm bảo tính mở và khả năng mở rộng linh hoạt \textit{[NFR5]} để sẵn sàng tích hợp thêm các dịch vụ giá trị gia tăng trong tương lai mà không cần tái cấu trúc nền tảng, giúp doanh nghiệp tối ưu hóa chi phí đầu tư dài hạn.